% ============================================================================
% Survival Skills (Test-Taking Tips) — 5 Practice Tests: Explorer's Journal
% Trail survival guide with compass-point navigation strategies
% Grade 7
% ============================================================================

\thispagestyle{empty}

{
\sffamily

% --- Title ---
\begin{center}
    \textcolor{funTealDark}{\fontsize{26}{30}\selectfont\bfseries
    \faCompass~Survival Skills~\faCompass}

    \smallskip
    \textcolor{funGray}{\normalsize Essential navigation strategies for every trail}
\end{center}

\vspace{3mm}

% --- The Compass Method (unique 4-point strategy) ---
\begin{tcolorbox}[
    enhanced,
    colback=funCream,
    colframe=funTealDark!50,
    arc=12pt, boxrule=2pt,
    left=5mm, right=5mm, top=4mm, bottom=4mm,
    title={\textcolor{white}{\large\faCompass~The Four-Point Compass}},
    colbacktitle=funTealDark!80,
    fonttitle=\bfseries\sffamily,
]
\small
\begin{center}
\begin{tikzpicture}[
    every node/.style={font=\sffamily\small, align=center},
    point/.style={
        draw=funTeal!40, fill=funTeal!8, thick,
        minimum width=5.5cm, minimum height=2.2cm,
        rounded corners=8pt, inner sep=3mm,
    },
]
    % North
    \node[point] (N) at (0, 2.8) {
        \textcolor{funTealDark}{\bfseries\faArrowUp~N --- Note the key info}\\
        Circle numbers. Underline the question.
    };
    % East
    \node[point] (E) at (5.5, 0) {
        \textcolor{funOrange}{\bfseries\faArrowRight~E --- Estimate first}\\
        Get a ballpark answer before solving.
    };
    % South
    \node[point] (S) at (0, -2.8) {
        \textcolor{funGreen}{\bfseries\faArrowDown~S --- Solve step by step}\\
        Write equations. Show every operation.
    };
    % West
    \node[point] (W) at (-5.5, 0) {
        \textcolor{funPurple}{\bfseries\faArrowLeft~W --- Weigh your answer}\\
        Does it match the question and estimate?
    };
    % Center compass icon
    \node[font=\fontsize{20}{20}\selectfont, funTealDark] at (0,0) {\faCompass};
\end{tikzpicture}
\end{center}
\end{tcolorbox}

\vspace{2mm}

% --- Grade 7-Specific Terrain Warnings ---
\begin{tcolorbox}[
    enhanced,
    colback=funYellowLight,
    colframe=funYellowDark!45,
    arc=10pt, boxrule=1pt,
    left=5mm, right=5mm, top=3mm, bottom=3mm,
    title={\small\textcolor{white}{\faExclamationTriangle~Grade 7 Terrain Warnings}},
    colbacktitle=funYellowDark!60,
    fonttitle=\bfseries\sffamily,
]
    \small
    \textcolor{funYellowDark}{\faExclamationTriangle}~\textbf{Negative numbers:} Watch signs when multiplying and dividing rational numbers.\\[2pt]
    \textcolor{funYellowDark}{\faExclamationTriangle}~\textbf{Proportions:} Make sure both ratios compare the \textit{same} quantities in the \textit{same} order.\\[2pt]
    \textcolor{funYellowDark}{\faExclamationTriangle}~\textbf{Percent problems:} Identify whether you need the part, the whole, or the percent.\\[2pt]
    \textcolor{funYellowDark}{\faExclamationTriangle}~\textbf{Geometry:} Label all measurements and don't confuse area with perimeter, or surface area with volume.\\[2pt]
    \textcolor{funYellowDark}{\faExclamationTriangle}~\textbf{Probability:} Read whether the question asks for theoretical or experimental probability.
\end{tcolorbox}

\vspace{2mm}

% --- Two-column: Timing + After the Trail ---
\noindent
\begin{minipage}[t]{0.48\linewidth}
\begin{tcolorbox}[
    enhanced,
    colback=funTeal!6,
    colframe=funTeal!35,
    arc=8pt, boxrule=0.8pt,
    left=4mm, right=4mm, top=3mm, bottom=3mm,
    title={\small\textcolor{white}{\faClock~Pacing Your Trek}},
    colbacktitle=funTeal!55,
    fonttitle=\bfseries\sffamily,
    height=5cm,
    valign=center,
]
    \small
    \begin{itemize}[leftmargin=*, itemsep=2pt]
        \item Trails 1--2: Untimed (explore freely)
        \item Trails 3--4: Timed (65 minutes)
        \item Trail 5: Strict test conditions
        \item Don't sprint --- steady pace wins
    \end{itemize}
\end{tcolorbox}
\end{minipage}%
\hfill
\begin{minipage}[t]{0.48\linewidth}
\begin{tcolorbox}[
    enhanced,
    colback=funOrange!6,
    colframe=funOrange!35,
    arc=8pt, boxrule=0.8pt,
    left=4mm, right=4mm, top=3mm, bottom=3mm,
    title={\small\textcolor{white}{\faFlag~After Each Trail}},
    colbacktitle=funOrange!55,
    fonttitle=\bfseries\sffamily,
    height=5cm,
    valign=center,
]
    \small
    \begin{itemize}[leftmargin=*, itemsep=2pt]
        \item Read every answer explanation
        \item Mark the \textit{topic}, not just the question
        \item Log your results in the Expedition Log
        \item Study rough terrain before the next trail
    \end{itemize}
\end{tcolorbox}
\end{minipage}

\vspace{3mm}

% --- Field note from guide ---
\begin{tcolorbox}[
    enhanced,
    colback=funCream,
    colframe=funTeal!25,
    arc=12pt, boxrule=0.8pt,
    left=5mm, right=5mm, top=3mm, bottom=3mm,
]
    \begin{minipage}[c]{0.07\linewidth}
        \centering \iconOwl[1cm]
    \end{minipage}%
    \hfill
    \begin{minipage}[c]{0.90\linewidth}
        \small\itshape\color{funTealDark}
        \faQuoteLeft~The best explorers don't avoid rough terrain ---
        they study it. Every wrong answer is a trailhead to
        a deeper understanding. Follow those paths.~\faQuoteRight
    \end{minipage}
\end{tcolorbox}

}

\newpage
